% Figura: verificación del teorema de la divergencia en una esfera
\begin{figure}[ht]
\centering
\begin{minipage}[t]{0.52\textwidth}\centering
\tdplotsetmaincoords{70}{125}
\begin{tikzpicture}[tdplot_main_coords, scale=1.05, baseline=(current bounding box.north)]
  \def\RR{2.1}
  \draw[eje] (0,0,0) -- (3.0,0,0) node[below left] {\small $x$};
  \draw[eje] (0,0,0) -- (0,3.1,0) node[right] {\small $y$};
  \draw[eje] (0,0,0) -- (0,0,3.1) node[above] {\small $z$};
  % esfera: contorno aparente y círculos máximos
  \begin{scope}[tdplot_screen_coords]
    \shade[ball color=azulclaro!55, opacity=0.35] (0,0) circle (\RR);
    \draw[azulprof!60, line width=0.6pt] (0,0) circle (\RR);
  \end{scope}
  \draw[azulprof!45, line width=0.4pt, domain=0:360, samples=72] plot ({\RR*cos(\x)},{\RR*sin(\x)},0);
  \foreach \ph in {0,45,90,135}{\draw[azulprof!35, line width=0.4pt, domain=0:360, samples=72]
     plot ({\RR*sin(\x)*cos(\ph)},{\RR*sin(\x)*sin(\ph)},{\RR*cos(\x)});}
  % campo radial A = R^2 u_r (flechas salientes, módulo constante)
  \foreach \th/\ph in {35/20,35/200,70/110,70/290,110/20,110/200,145/110,145/290}{%
    \pgfmathsetmacro{\px}{\RR*sin(\th)*cos(\ph)}\pgfmathsetmacro{\py}{\RR*sin(\th)*sin(\ph)}\pgfmathsetmacro{\pz}{\RR*cos(\th)}%
    \draw[-{Stealth[length=4pt]}, line width=1pt, acento] (\px,\py,\pz) -- ++({0.55*sin(\th)*cos(\ph)},{0.55*sin(\th)*sin(\ph)},{0.55*cos(\th)});}
  % elemento de superficie y normal
  \def\ta{48}\def\tb{62}\def\pa{28}\def\pb{46}
  \fill[ocreteo, opacity=0.7]
    ({\RR*sin(\ta)*cos(\pa)},{\RR*sin(\ta)*sin(\pa)},{\RR*cos(\ta)}) --
    ({\RR*sin(\tb)*cos(\pa)},{\RR*sin(\tb)*sin(\pa)},{\RR*cos(\tb)}) --
    ({\RR*sin(\tb)*cos(\pb)},{\RR*sin(\tb)*sin(\pb)},{\RR*cos(\tb)}) --
    ({\RR*sin(\ta)*cos(\pb)},{\RR*sin(\ta)*sin(\pb)},{\RR*cos(\ta)}) -- cycle;
  \draw[-{Stealth[length=5pt]}, line width=1.5pt, azulacero]
    ({\RR*sin(55)*cos(37)},{\RR*sin(55)*sin(37)},{\RR*cos(55)}) --
    ({(\RR+1.0)*sin(55)*cos(37)},{(\RR+1.0)*sin(55)*sin(37)},{(\RR+1.0)*cos(55)});
  \node[font=\small, text=azulacero, anchor=south west] at ({(\RR+1.0)*sin(55)*cos(37)},{(\RR+1.0)*sin(55)*sin(37)},{(\RR+1.0)*cos(55)}) {$\dS$};
  \fill (0,0,0) circle (1.2pt);
\end{tikzpicture}\\[2pt]{\small\scshape (a)}
\end{minipage}\hfill
\begin{minipage}[t]{0.44\textwidth}\centering
\begin{tikzpicture}[baseline=(current bounding box.north)]
  \begin{axis}[tema, width=\textwidth, height=6.3cm, xlabel={$r$}, xmin=0, xmax=1.6, ymin=0, ymax=42,
      samples=100, domain=0:1.6, xtick={0,0.5,1,1.5}, xticklabels={0,0{,}5,$R$,1{,}5},
      xtick={0,0.5,1.0,1.5}, legend style={at={(0.03,0.97)}, anchor=north west}]
    \addplot[azulacero, line width=1.1pt] {4*pi*x^4};
    \addplot[acento, line width=1.1pt] {16*pi*x^3};
    \addplot[name path=curva, acento, draw=none, domain=0:1] {16*pi*x^3};
    \path[name path=eje] (axis cs:0,0) -- (axis cs:1,0);
    \addplot[acento!25] fill between[of=curva and eje];
    \addplot[grispizarra, densely dashed] coordinates {(1,0) (1,12.57)};
    \node[font=\scriptsize, text=azulacero, anchor=west] at (axis cs:1.05,9.5) {$4\pi R^{4}$};
    \legend{$\Phi(r)=4\pi r^{4}$, $16\pi r^{3}=\nab\cdot\vc{A}\,\cdot4\pi r^{2}$}
  \end{axis}
\end{tikzpicture}\\[2pt]{\small\scshape (b)}
\end{minipage}
\caption{Teorema de Gauss para $\vc{A}=r^{2}\,\ur$ en una esfera de radio
$R$. (a)~El campo es radial y de módulo constante sobre la esfera; la normal
exterior es $\ur$, de modo que $\vc{A}\cdot\un=R^{2}$ y el flujo es
$R^{2}\cdot4\pi R^{2}$. (b)~Visión «capa a capa»: cada cáscara de radio $r$ y
espesor $\dif r$ aporta $\nab\cdot\vc{A}\,\dV=4r\cdot4\pi r^{2}\,\dif r$; el área
bajo $16\pi r^{3}$ (sombreada) hasta $R=1$ es el flujo total
$\Phi(R)=4\pi R^{4}$.}
\label{fig:sem-esfera}
\end{figure}
